% ============================================================
%  Aircraft Fleet Predictive Maintenance System
%  Engineering Technical Report -- Revision A
%  Author: Robert A. -- Aerospace Engineering Graduate
% ============================================================
\documentclass[11pt, a4paper]{report}

% --- Packages -----------------------------------------------
\usepackage[margin=2.5cm]{geometry}
\usepackage{amsmath, amssymb}
\usepackage{booktabs}
\usepackage{array}
\usepackage{longtable}
\usepackage{graphicx}
\usepackage{hyperref}
\usepackage{xcolor}
\usepackage{enumitem}
\usepackage{fancyhdr}
\usepackage{titlesec}
\usepackage{caption}
\usepackage{multirow}
\usepackage{microtype}
\usepackage{parskip}

% --- Colours ------------------------------------------------
\definecolor{dimgray}{RGB}{74, 90, 122}

% --- Header / Footer ----------------------------------------
\pagestyle{fancy}
\fancyhf{}
\fancyhead[L]{\small\texttt{SKYPULSE AI -- ENGINEERING TECHNICAL REPORT}}
\fancyhead[R]{\small\texttt{REV A / JUNE 2026}}
\fancyfoot[C]{\small\thepage}
\renewcommand{\headrulewidth}{0.4pt}

% --- Section formatting -------------------------------------
\titleformat{\chapter}[block]
  {\normalfont\Large\bfseries}{}{0em}
  {\thechapter\quad}[\vspace{0.5ex}\hrule\vspace{1ex}]
\titleformat{\section}{\normalfont\large\bfseries}{\thesection}{1em}{}
\titleformat{\subsection}{\normalfont\normalsize\bfseries}{\thesubsection}{1em}{}

% --- Hyperref -----------------------------------------------
\hypersetup{
  colorlinks=true,
  linkcolor=black,
  urlcolor=blue,
  citecolor=black
}

% ============================================================
\begin{document}

% --- Title page ---------------------------------------------
\begin{titlepage}
  \centering
  \vspace*{2cm}
  {\small\texttt{TECHNICAL REPORT / ENGINEERING PROOF-OF-CONCEPT}}\\[0.4cm]
  \hrule height 0.6pt
  \vspace{1cm}
  {\LARGE\bfseries Aircraft Fleet Predictive Maintenance System:\\[0.3cm]
  A Reliability-Engineering-Based Proof of Concept}\\[1.5cm]
  \hrule height 0.6pt
  \vspace{1.5cm}
  \begin{tabular}{ll}
    \textbf{Author}  & Robert A. -- Aerospace Engineering Graduate \\[4pt]
    \textbf{Date}    & June 2026 \\[4pt]
    \textbf{Revision}& A \\[4pt]
    \textbf{Status}  & Educational Proof-of-Concept -- Not for Operational Use \\
  \end{tabular}
  \vspace{2cm}

  \begin{tabular}{|l|l|}
    \hline
    \textbf{Parameter}  & \textbf{Value} \\ \hline
    Aircraft simulated  & 8 \\ \hline
    Subsystems per A/C  & 6 \\ \hline
    Engineering models  & 5 (health, RUL, reliability, VSI, FMEA) \\ \hline
    Sensor parameters   & 10 per aircraft \\ \hline
    Telemetry window    & 48 readings, 4-second ticks \\ \hline
  \end{tabular}

  \vfill
  {\small\color{dimgray}
    \textbf{DISCLAIMER:} This is a software simulation using procedurally generated data.
    No results should be used for real aircraft maintenance or airworthiness decisions.
    Operational PHM systems require certified data acquisition, validated degradation models,
    and compliance with FAA AC 120-66B / EASA Part-M.
  }
\end{titlepage}

\tableofcontents
\newpage

% ============================================================
\chapter{Executive Summary}

This report documents a proof-of-concept PHM (Prognostics and Health Management)
system for aircraft fleet maintenance. The system takes simulated sensor data from
eight aircraft, computes subsystem health scores, estimates Remaining Useful Life
(RUL) using a stress-adjusted model, calculates mission reliability via the
exponential function, and ranks maintenance tasks by FMEA risk score.

The focus throughout is on keeping the calculations explicit and traceable. Every
number on the dashboard can be derived by hand from the formulas in
Chapter~\ref{ch:methodology}, and the worked example in Chapter~\ref{ch:worked}
does exactly that for a single aircraft.

\medskip
\noindent\textbf{Main finding:} Rule-based PHM with explicit engineering models
gives interpretable, auditable results. The gap to a production system is real
but well-defined: it requires actual sensor data, validated degradation curves,
and regulatory approval. That gap is documented in Chapters~\ref{ch:limitations}
and~\ref{ch:production}.

% ============================================================
\chapter{Problem Statement}

Unscheduled maintenance events cost the aviation industry roughly \$9 billion
per year in direct costs and disruption (IATA, 2023). Most of this comes from
two problems with the standard scheduled-maintenance approach:

\begin{enumerate}
  \item \textbf{Over-maintenance:} Parts get replaced before they need to be,
        which wastes money and maintenance hours.
  \item \textbf{Under-maintenance:} Parts that degrade faster than the rated
        interval fail before the next scheduled check.
\end{enumerate}

Condition-Based Maintenance (CBM) fixes both by estimating actual component
health from sensor data and scheduling work based on real condition rather than
calendar time. This project builds a simplified version of that calculation
chain to show how the underlying engineering works.

% ============================================================
\chapter{Scope and Objectives}

\section{What This Project Covers}

\begin{itemize}
  \item Simulated telemetry for 8 aircraft across 4 health scenarios
        (nominal, degraded, warning, critical).
  \item Health score computation for 6 subsystems per aircraft.
  \item Stress-adjusted RUL estimation for the engine.
  \item Mission reliability using the exponential model.
  \item FMEA RPN scoring for maintenance prioritisation.
  \item A live web interface showing all computed values with inline
        calculation breakdowns.
\end{itemize}

\section{What This Project Does Not Cover}

\begin{itemize}
  \item Real sensor data or aircraft bus integration.
  \item Calibrated degradation models from fleet history.
  \item Regulatory compliance (FAA, EASA, MIL-HDBK-217).
  \item Weibull fitting or Bayesian parameter estimation.
  \item Maintenance cost modelling or scheduling optimisation.
\end{itemize}

% ============================================================
\chapter{Assumptions and Justifications}

\begin{longtable}{@{}p{5cm} p{4.5cm} p{5cm}@{}}
  \toprule
  \textbf{Assumption} & \textbf{Value / Form} & \textbf{Justification} \\
  \midrule
  \endhead
  Constant failure rate &
  $R(t) = e^{-\lambda t}$ &
  Valid in the useful-life phase of the bathtub curve. Simplest tractable model;
  standard in early-phase reliability work. \\
  \addlinespace
  Linear interpolation between sensor thresholds &
  $H = 80 + \frac{(x_{warn} - x)}{(x_{warn} - x_{nom})} \times 20$ &
  No calibrated degradation data available. Linear is conservative and easy to
  audit. \\
  \addlinespace
  Scalar stress factor from temperature and vibration &
  $SF = \frac{T_{factor} + V_{factor}}{2}$ &
  Simplified proxy for Miner's rule (fatigue) and Arrhenius (thermal). Both
  effects are known to accelerate component aging above nominal conditions. \\
  \addlinespace
  Temperature stress coefficient $k_T = 0.5$ &
  $T_{factor} = 1 + (1 - T_{score}) \times 0.5$ &
  Represents Arrhenius thermal activation. The specific coefficient is
  illustrative; real values are alloy- and failure-mode-specific. \\
  \addlinespace
  Vibration stress coefficient $k_V = 0.8$ &
  $V_{factor} = 1 + (1 - V_{score}) \times 0.8$ &
  Set higher than temperature because vibration drives mechanical fatigue via
  Miner's rule more directly at sub-critical amplitudes. \\
  \addlinespace
  Engine rated life = 25,000 FH &
  OEM life limit &
  Typical TBO for a CFM56-class turbofan. Used as a reference only. \\
  \addlinespace
  MTBF from point estimate &
  $\text{MTBF} = T_{total} / N_{failures}$ &
  Standard MLE estimator for the exponential distribution. Known limitation:
  wide CI at small $N$ (see Section~\ref{sec:mtbf}). \\
  \addlinespace
  VSI warning at 2.5 &
  $\text{VSI}_{warn} = 2.5$ &
  ISO 10816-3 Zone B extends to 2.8 mm/s. Using 2.5 as an early-warning margin. \\
  \addlinespace
  FMEA severity fixed per subsystem &
  Engine: $S = 9$ &
  Simplified. A real FMEA assigns severity per failure mode, not per subsystem. \\
  \bottomrule
\end{longtable}

% ============================================================
\chapter{Engineering Methodology}
\label{ch:methodology}

\section{Sensor Normalisation}

Each sensor reading is mapped to a $[0, 100]$ score using piecewise linear
interpolation across three threshold values: nominal ($x_{nom}$), warning
($x_{warn}$), and critical ($x_{crit}$). Scores above 80 are in the nominal
band; below 40 is critical.

\subsection{Scoring Function}

\[
H(x) =
\begin{cases}
100 & x \leq x_{nom} \\[4pt]
80 + 20 \cdot \dfrac{x_{warn} - x}{x_{warn} - x_{nom}} & x_{nom} < x \leq x_{warn} \\[8pt]
40 + 40 \cdot \dfrac{x_{crit} - x}{x_{crit} - x_{warn}} & x_{warn} < x \leq x_{crit} \\[8pt]
\max\!\left(0,\; 40 \cdot \dfrac{x}{x_{crit}}\right) & x > x_{crit}
\end{cases}
\]

For sensors where lower is worse (oil pressure, battery voltage), the
inequalities are reversed.

\subsection{Sensor Thresholds}

\begin{table}[h!]
\centering
\begin{tabular}{@{}llrrrl@{}}
\toprule
\textbf{Subsystem} & \textbf{Sensor} & \textbf{Nominal} & \textbf{Warning} & \textbf{Critical} & \textbf{Units} \\
\midrule
Engine      & Temperature    & 650  & 750  & 850  & $^\circ$C \\
            & Vibration      & 2.0  & 5.0  & 10.0 & mm/s \\
            & Oil Pressure   & 60   & 45   & 30   & PSI \\
Hydraulic   & Pressure       & 3000 & 2400 & 1800 & PSI \\
Electrical  & Battery Voltage& 28.0 & 25.5 & 23.0 & V \\
            & Generator Volt & 115  & 112  & 108  & V AC \\
Avionics    & Temperature    & 45   & 60   & 75   & $^\circ$C \\
ECS         & Temperature    & 45   & 58   & 72   & $^\circ$C \\
Landing Gear& Cycles         & 0    & --   & 50000& cycles \\
\bottomrule
\end{tabular}
\caption{Sensor thresholds by subsystem.}
\end{table}

\section{Subsystem Health Scores}

Each subsystem score is a weighted sum of sensor scores. Weights reflect the
relative diagnostic importance of each sensor for that subsystem.

\subsection{Engine}

\[
H_{engine} = 0.30 \cdot h_T + 0.30 \cdot h_V + 0.20 \cdot h_P + 0.20 \cdot h_U
\]

Temperature and vibration share equal top weight at 0.30 each because both
are direct indicators of thermomechanical degradation in turbofans.
$h_U$ is the usage fraction score, where $h_U = 100 \cdot (1 - FH/L_{rated})$.

\subsection{Other Subsystems}

\begin{align}
  H_{hydraulic}   &= 0.60 \cdot h_P + 0.40 \cdot h_U \\[4pt]
  H_{electrical}  &= 0.40 \cdot h_{bat} + 0.40 \cdot h_{gen} + 0.20 \cdot h_U \\[4pt]
  H_{avionics}    &= 0.60 \cdot h_T + 0.40 \cdot h_U \\[4pt]
  H_{ECS}         &= 0.50 \cdot h_T + 0.50 \cdot h_U \\[4pt]
  H_{gear}        &= 0.70 \cdot h_{cyc} + 0.30 \cdot h_U
\end{align}

\section{Vibration Severity Index}

Per ISO 10816:

\[
\text{VSI} = \frac{V_{measured}}{V_{nominal}} = \frac{V_{measured}}{2.0 \text{ mm/s}}
\]

\begin{table}[h!]
\centering
\begin{tabular}{@{}clll@{}}
\toprule
\textbf{Zone} & \textbf{VSI Range} & \textbf{Classification} & \textbf{Action} \\
\midrule
A & $\leq 1.8$   & Good                        & None \\
B & $1.8 - 2.8$  & Acceptable                  & Monitor trend \\
C & $2.8 - 4.5$  & Unsatisfactory              & Investigate; schedule maintenance \\
D & $> 4.5$      & Unacceptable                & Remove from service \\
\bottomrule
\end{tabular}
\caption{ISO 10816 vibration severity zones.}
\end{table}

\section{Stress Factor and Remaining Useful Life}

\subsection{Stress Factor}

The stress factor converts actual flight hours into equivalent damage hours.
An engine running hot and rough accumulates fatigue faster than its rated
conditions assume.

\begin{align}
  T_{factor} &= 1.0 + (1 - h_T/100) \times 0.5 \\[4pt]
  V_{factor} &= 1.0 + (1 - h_V/100) \times 0.8 \\[4pt]
  SF         &= \frac{T_{factor} + V_{factor}}{2}
\end{align}

$SF = 1.0$ is nominal. $SF = 1.1$ means the component ages 10\% faster than
its life limit was written for.

\subsection{Remaining Useful Life}

\begin{align}
  \text{Effective Age} &= FH \times SF \\[4pt]
  \text{RUL}           &= L_{rated} - \text{Effective Age}
\end{align}

\subsection{Fatigue Fraction Consumed}

\[
\text{FAT} = \frac{FH \times SF}{L_{rated}} \times 100\%
\]

\section{MTBF and Mission Reliability}
\label{sec:mtbf}

\[
\text{MTBF} = \frac{T_{total}}{N_{failures}}, \quad
\lambda = \frac{1}{\text{MTBF}}, \quad
R(t) = e^{-\lambda t}
\]

$R(t)$ is the probability of completing a $t$-hour mission without failure.
The model assumes a constant hazard rate, which holds during the useful-life
phase only. As components approach their life limits the hazard rate increases,
and $R(t)$ will overestimate reliability.

With $N < 30$ failures the 90\% CI on MTBF spans roughly a factor of 5--7,
so the point estimate should not be used for safety-critical planning at
this sample size.

\section{FMEA Risk Priority Number}

\[
\text{RPN} = S \times P \times D
\]

\begin{itemize}
  \item $S$: Severity (1--10). Fixed per subsystem in this implementation.
  \item $P$: Probability (1--10), derived from health score:
        $P = \lceil (100 - H_{sub})/100 \times 10 \rceil$.
  \item $D$: Detectability (1--10). Assigned based on how visible the
        degradation is at the current health level.
\end{itemize}

\begin{table}[h!]
\centering
\begin{tabular}{@{}lll@{}}
\toprule
\textbf{Risk Tier} & \textbf{RPN Range} & \textbf{Response} \\
\midrule
CRITICAL & $\geq 600$  & Immediate grounding or maintenance action \\
HIGH     & $300 - 599$ & Maintenance within 50 FH \\
MEDIUM   & $100 - 299$ & Maintenance within 200 FH \\
LOW      & $< 100$     & Action at next scheduled interval \\
\bottomrule
\end{tabular}
\caption{RPN classification and required response.}
\end{table}

% ============================================================
\chapter{Worked Example -- Aircraft AC-102}
\label{ch:worked}

AC-102 (N7734W, Airbus A320-200) is a mid-life aircraft in a degraded scenario.
Input state:

\begin{table}[h!]
\centering
\begin{tabular}{@{}llll@{}}
\toprule
\textbf{Parameter}  & \textbf{Value} & \textbf{Parameter} & \textbf{Value} \\
\midrule
Flight Hours        & 14,200 hrs     & Battery Voltage    & 26.5 V \\
Flight Cycles       & 9,800          & Generator Voltage  & 113.2 V AC \\
Engine Temperature  & 720$^\circ$C   & Avionics Temp      & 57$^\circ$C \\
Engine Vibration    & 4.2 mm/s       & Landing Gear Cyc.  & 9,712 \\
Oil Pressure        & 51 PSI         & Recorded Failures  & 5 \\
Hydraulic Pressure  & 2,650 PSI      & Engine Life Limit  & 25,000 hrs \\
\bottomrule
\end{tabular}
\caption{AC-102 sensor snapshot used in all calculations below.}
\end{table}

\section{Engine Health Score}

\[
H_{engine} = 0.30 \cdot h_T + 0.30 \cdot h_V + 0.20 \cdot h_P + 0.20 \cdot h_U
\]

\paragraph{Temperature score.} $T = 720^\circ$C; nominal = 650, warning = 750.

\[
h_T = 80 + \frac{750 - 720}{750 - 650} \times 20 = 80 + 6.0 = 86
\]

\paragraph{Vibration score.} $V = 4.2$ mm/s; nominal = 2.0, warning = 5.0.

\[
h_V = 80 + \frac{5.0 - 4.2}{5.0 - 2.0} \times 20 = 80 + 5.3 = 85.3 \rightarrow 85
\]

\paragraph{Oil pressure score.} $P = 51$ PSI; nominal = 60, warning = 45. Higher is better.

\[
h_P = 80 + \frac{51 - 45}{60 - 45} \times 20 = 80 + 8.0 = 88
\]

\paragraph{Usage score.} $FH = 14{,}200$, $L_{rated} = 25{,}000$.
Usage fraction = 0.568, which is below 75\%, so the score stays in the
$[80, 100]$ band:

\[
h_U = 80 + 20 \times \frac{0.75 - 0.568}{0.75} = 80 + 4.9 = 84.9 \rightarrow 85
\]

\paragraph{Composite.}

\[
H_{engine} = 0.30(86) + 0.30(85) + 0.20(88) + 0.20(85)
           = 25.8 + 25.5 + 17.6 + 17.0 = \mathbf{85.9 \approx 86\%}
\]

\medskip
\noindent\textbf{Note:} 86\% is LOW risk and correct for the maintenance
schedule. What the composite score misses is that $h_V = 85$ and $h_P = 88$
are both close to their warning thresholds, and they are correlated. Simultaneous
vibration increase and oil pressure drop is consistent with bearing wear: the
degrading surface increases dynamic load while also letting oil bypass through
widened clearances. Neither sensor alone triggers a flag; together they are a
meaningful signal.

\section{Vibration Severity Index}

\[
\text{VSI} = \frac{4.2}{2.0} = \mathbf{2.10}
\]

Zone B (1.8--2.8): acceptable for unrestricted operation.
Margin to Zone C: $2.8 - 2.10 = 0.70$ mm/s.

\medskip
\noindent\textbf{Note:} The point value is fine. What matters is the rate.
If VSI has drifted from 1.8 to 2.10 over the last 300 FH (+0.1/100 hrs),
Zone C is roughly 700 FH away. I would set up a trend plot and define an
intervention threshold before that point rather than waiting for the threshold
to be crossed.

\section{Stress Factor and RUL}

\begin{align}
  T_{factor} &= 1.0 + (1 - 86/100) \times 0.5 = 1.07 \\[4pt]
  V_{factor} &= 1.0 + (1 - 85/100) \times 0.8 = 1.12 \\[4pt]
  SF         &= \frac{1.07 + 1.12}{2} = \mathbf{1.095}
\end{align}

\begin{align}
  \text{Effective Age} &= 14{,}200 \times 1.095 = 15{,}549 \text{ hrs} \\[4pt]
  \text{RUL}           &= 25{,}000 - 15{,}549 = \mathbf{9{,}451 \text{ hrs}}
\end{align}

Without the stress correction: $25{,}000 - 14{,}200 = 10{,}800$ hrs.
The stress factor removes 1,349 hrs. At 2,700 FH/year utilisation, that is
about six months of additional life consumed compared to what a naive calculation
would show.

\section{MTBF and Mission Reliability}

\begin{align}
  \text{MTBF} &= \frac{14{,}200}{5} = \mathbf{2{,}840 \text{ hrs}} \\[4pt]
  \lambda     &= \frac{1}{2{,}840} = 3.52 \times 10^{-4} \text{ failures/hr} \\[4pt]
  R(100)      &= e^{-0.000352 \times 100} = e^{-0.0352} = \mathbf{96.5\%}
\end{align}

\medskip
\noindent\textbf{Note:} With only 5 failures, the 90\% CI on MTBF spans roughly
1,200--8,500 hrs. The 2,840-hr figure is the best available estimate, but I
would not use it for safety-critical planning. Typically 30--50 events are
needed before MTBF feeds into RCM decisions.

$R(100) = 96.5\%$ holds under constant-failure-rate. As AC-102 approaches
its 25,000-hr life limit, wear-out will drive $\lambda$ higher and the true
reliability will be lower than this.

\section{FMEA RPN -- Engine}

\begin{align}
  S   &= 9 \quad \text{(safety-critical failure mode)} \\
  P   &= \lceil (100 - 86) / 100 \times 10 \rceil = \lceil 1.4 \rceil = 2 \\
  D   &= 5 \quad \text{(degradation visible but not alarming)} \\[4pt]
  \text{RPN} &= 9 \times 2 \times 5 = \mathbf{90 \text{ -- LOW}}
\end{align}

\medskip
\noindent\textbf{Note:} RPN = 90 is correct for formal scheduling. The issue
is that RPN scores the current snapshot, not the direction of travel. An RPN of
90 on a stable aircraft is different from the same score on one with rising
vibration and dropping oil pressure. I would keep the LOW classification but add
a monitoring flag for the correlated sensor pattern.

\section{Landing Gear Fatigue Life}

\[
\text{FAT} = \frac{9{,}712}{50{,}000} \times 100 = \mathbf{19.4\%}
\]

At roughly 1,225 cycles/year, the structural life limit is 33 years away.
For AC-102, the engine RUL (9,451 hrs) is the binding constraint, not the
airframe.

% ============================================================
\chapter{Engineering Assessment -- AC-102}

Engine Health 86\% puts AC-102 in the LOW risk category. That is the correct
formal classification.

The issue is what the composite score does not show. Vibration is at VSI 2.10
(0.70 mm/s from Zone C) and oil pressure is at 51 PSI (6 PSI above the warning
threshold). Both are nominally acceptable. The combination is not a coincidence:
these two parameters moving in opposite directions simultaneously is the
fingerprint of early bearing wear. The weighted average gets pulled up by the
four healthy sensors and returns 86\%, which could lead a technician to miss it.

\medskip
\noindent\textbf{Recommended actions:}

\begin{enumerate}
  \item Keep LOW classification for the maintenance schedule.
  \item Add to the next C-check borescope list.
  \item Set up a trend plot for vibration and oil pressure.
        Intervention thresholds: VSI 2.5, oil pressure 48 PSI.
  \item If either crosses its warning threshold before the next scheduled
        check, escalate to HIGH and action within 50 FH.
\end{enumerate}

\medskip
\noindent Primary life limit: Engine RUL = 9,451 FH (stress-adjusted). \\
Non-binding: Landing gear (19.4\%), electrical, hydraulic -- all nominal.

% ============================================================
\chapter{Model Validation and Sanity Checks}

\section{Automated Tests}

18 unit tests cover all calculation modules. These check properties, not
empirical accuracy:

\begin{itemize}
  \item Health scores bounded to $[0, 100]$.
  \item $R(t) \in (0, 1]$ for all $t \geq 0$, $R(t) \to 0$ as $t \to \infty$.
  \item Stress factor $SF \geq 1.0$ for all valid sensor readings.
  \item RUL floored at 0.
  \item RPN in $[1, 1000]$.
\end{itemize}

\section{Dimensional Consistency}

\begin{itemize}
  \item Health scores: dimensionless, $[0, 100]$.
  \item $\lambda$: hr$^{-1}$. $\lambda t$ is dimensionless. Correct.
  \item RUL: flight hours. $FH \times SF$ is hours, $L_{rated}$ is hours. Correct.
  \item VSI: dimensionless (mm/s divided by mm/s). Correct.
\end{itemize}

\section{Boundary Behaviour}

\begin{itemize}
  \item $FH = 0$: usage score = 100, $SF = 1.0$, RUL = $L_{rated}$. Correct.
  \item $FH = L_{rated}$: RUL = 0. Correct.
  \item All sensors at nominal: all scores = 100, $SF = 1.0$. Correct.
  \item $t = 0$: $R(0) = 1.0$. Correct.
\end{itemize}

% ============================================================
\chapter{Limitations}
\label{ch:limitations}

\begin{enumerate}
  \item \textbf{No real sensor data.} Telemetry is procedurally generated from
        seeded pseudo-random functions. No model parameters are calibrated to a
        real aircraft type or failure dataset.

  \item \textbf{Exponential reliability model.} Constant hazard rate is only
        valid in the useful-life phase. Late-life wear-out is better described
        by a Weibull distribution ($\beta > 1$). The current model will
        overestimate reliability as components approach their life limits.

  \item \textbf{Scalar stress factor.} This is a simplified approximation. A
        production system would use physics-of-failure models (Paris law for
        crack growth, Arrhenius for thermal degradation) calibrated to material
        data and fleet history.

  \item \textbf{MTBF confidence.} With $N < 10$ failures per aircraft, confidence
        intervals are wide. The uncertainty is not propagated to downstream
        calculations; the point estimate is used throughout.

  \item \textbf{Fixed sensor weights.} Weights are engineering judgments, not
        regression-derived. A production system would determine weights from
        failure outcome data or formal importance analysis.

  \item \textbf{No cross-correlation between sensors.} Correlated failure
        signatures (like the AC-102 vibration + oil pressure case) are not
        explicitly detected. The weighted average can mask them.

  \item \textbf{Not regulatory compliant.} This does not meet FAA AC 120-66B,
        EASA Part-M, or MSG-3 requirements for maintenance decision support.
\end{enumerate}

% ============================================================
\chapter{Steps to a Production System}
\label{ch:production}

\begin{enumerate}
  \item \textbf{Real data acquisition.} Interface with ARINC 429 / ARINC 664
        or QAR download. Establish data traceability and timestamp sync.

  \item \textbf{Weibull fitting.} Replace the exponential model with a Weibull
        fit to fleet failure history. Estimate $\beta$ and $\eta$ by MLE.
        Update with each new failure event using Bayesian methods.

  \item \textbf{Degradation model calibration.} Replace linear interpolation
        with physics-of-failure curves validated against real failure events.

  \item \textbf{Uncertainty propagation.} All RUL estimates should carry
        confidence bounds, not just point values. Monte Carlo propagation from
        sensor uncertainty to RUL distribution is the minimum standard.

  \item \textbf{Formal FMEA.} Conduct FMEA/FMECA per ARP 4761. Assign $S$,
        $P$, $D$ per failure mode rather than per subsystem.

  \item \textbf{Regulatory approval.} Submit to FAA/EASA for health management
        system approval. Comply with DO-178C (software) and DO-254 (hardware).

  \item \textbf{MRO integration.} Connect to the airline's MMS for automatic
        work order generation and parts provisioning.
\end{enumerate}

% ============================================================
\chapter{Conclusions}

The calculation chain from raw sensor reading to maintenance recommendation is
fully explicit and reproducible. Every output in the web interface can be
derived by hand from the formulas in Chapter~\ref{ch:methodology}, and the
worked example in Chapter~\ref{ch:worked} does that for AC-102 in full.

The AC-102 analysis surfaced a practical limitation of composite health scoring:
the overall score of 86\% is technically correct, but it obscures the correlated
vibration and oil pressure trend that a maintainer should actually be watching.
This is not unique to this implementation. Any weighted-average health score
has this property. The fix is to display the individual sensor inputs alongside
the composite, which the interface does.

The gap between this proof-of-concept and a deployable system is real and
documented in Chapter~\ref{ch:limitations}. The main items are real sensor data,
calibrated degradation models, and regulatory approval. None of those are
unsolvable engineering problems; they just require resources and fleet history
that are outside the scope of a portfolio project.

\vfill
\noindent\hrule\vspace{0.5em}
\small
\noindent\texttt{END OF REPORT -- REV A}
\hfill
\texttt{Aircraft Fleet Predictive Maintenance System / Robert A. / June 2026}

\end{document}
